\subsection{STM32F405}
\label{chap:analysis_stm32f05}


The STM32F405[ref], part of STMicroelectronics' STM32F4 family, is based 
on a 32-bit ARM Cortex-M4 core with a Floating Point Unit (FPU), support for 
DSP instructions, an Adaptive Real-Time Accelerator (ART) for 0-wait execution 
from Flash up to 168 MHz, a Memory Protection Unit (MPU), and 210 DMIPS 
of performance. This configuration makes it an ideal microcontroller for embedded 
systems with real-time and intensive processing requirements, such as motor control, 
power conversion, and multimedia applications. It incorporates up to 1 MB of Flash 
memory, up to 192+4 KB of SRAM (with 64 KB of CCM), and essential peripherals: 
three 12-bit ADCs, two DACs, 17 timers, 16-stream DMA, three I$^2$C ports, multiple 
USART/SPI/CAN ports, Ethernet MAC, USB OTG HS/FS, and up to 140 fast, 
5V-tolerant GPIOs.

\subsubsection{STM32F405 I$^2$C}

As explained previously, It has been decided to add emulation of the I$^2$C \\
STM32F405 communication protocol to QEMU. This microcontroller has three I$^2$C 
peripherals (I$^2$C1, I$^2$C2, I$^2$C3), each providing a fully compliant I$^2$C bus interface. 
These peripherals automatically manage bus-specific sequencing, protocol 
management, arbitration, and timing, simplifying firmware development and ensuring 
robust communication. A Summary of the main capabilities of the STM32F405\cite{STM32F405Datasheet2026}:

\begin{itemize}
    \item I$^2$C standard compliance, with 7-bit and 10-bit addressing modes.
    \item Dual-speed modes: 
        \begin{itemize}
            \item Standard mode(Sm), up to 100 kHz.
            \item Fast mode(Fm) up to 400 kHz.
        \end{itemize}
    \item Master-slave operation, with multi-master capability and clock expansion support.
    \item Hardware packet error checking (PEC) using the CRC-8 polynomial for increased reliability.
    \item DMA support for efficient data transfers without CPU intervention. 
    \item Dual interrupt vectors: one for successful communication events and one for error conditions.
\end{itemize}

A design decision made prior to implementation was to focus I$^2$C development 
exclusively on master mode, as this represents the operating context in which the 
STM32F405 communicates with the LIS3DH accelerometer. The emulated 
microcontroller's firmware initiates all I$^2$C transactions, while the LIS3DH device model 
responds as a slave. This design decision simplifies the scope of the emulation while 
maintaining full functional reliability for the intended use case.


\subsubsection{STM32F405 I$^2$C Registers}

The STM32F405 I$^2$C peripheral operates through a set of memory-mapped 
registers that control all aspects of bus communication\cite{STM32F405ReferenceManual2021}. Understanding this register 
architecture is fundamental to accurate emulation, as firmware interacts with the 
peripheral exclusively through register read and write operations. The register set 
can be functionally categorized into four groups:

\begin{itemize}
    \item \textbf{Control Registers:} The registers I2C\_CR1 and I2C\_CR2
    \item \textbf{Timing Configuration Registers:} The registers I2C\_CCR and I2C\_TRISE
    \item \textbf{Data and Address Registers:} The registers I2C\_DR and I2C\_OAR1
    \item \textbf{Status Registers:} The registers I2C\_SR1 and I2C\_SR2
\end{itemize}

Among these register groups, the control registers, status registers, and the 
I2C\_DR register will have the greatest impact on emulation. On the other hand, 
the I2C\_OAR1 register is irrelevant, as it only defines the slave address of the 
peripheral when the microcontroller is operating in slave mode. Regarding the timing 
configuration registers, it was decided not to implement their associated logic. This 
is because QEMU, to emulate the operation of the I$^2$C protocol, provides a series 
of functions that execute the different stages of the I$^2$C communication protocol 
without considering timing. Therefore, it is not currently necessary to analyze 
these registers.

For implementation in the emulator, it is crucial to understand the logic of the 
control and status registers, as well as the I2C\_DR register. The I2C\_DR register, 
the data register, functions as a bidirectional data buffer between the firmware and the 
I$^2$C shift register. The control registers, I2C\_CR1 and I2C\_CR2, on the other hand, 
serve as the main interface for peripheral control and trigger transitions between 
the different phases of the I$^2$C protocol. Specifically, the I2C\_CR2 register contains 
the bits that configure clock generation and interrupt control, while the I2C\_CR1 
register contains, among others, the following bits:

\begin{itemize}
    \item \textbf{PE}, which enables the use of I$^2$C communication itself.
    \item \textbf{START}, which generates the start condition on the bus.
    \item \textbf{STOP}, which is responsible for generating the stop condition on the bus and terminating communication.
    \item \textbf{ACK}, the acknowledgment bit, which controls whether the peripheral generates an acknowledgment pulse after receiving a byte, which is essential for multibyte read operations.
\end{itemize}


Following the previous explanation, the status registers contain information about 
certain conditions that have been met during a specific operation and that can be 
taken into account in subsequent operations. The I2C\_SR2 register contains bits 
with information about the bus status, including the MSL bit, which indicates the 
current operating mode, the BUSY bit, which reflects bus activity, and the TRA 
bit, which shows the current data direction. The I2C\_SR1 register, on the other 
hand, provides real-time visibility into communication events and error conditions. 
The key indicators in this register are:

\begin{itemize}
    \item \textbf{SB:     }Set when a START condition (master mode) has been generated.
    \item \textbf{ADDR:   }Set when the address byte has been successfully transmitted.
    \item \textbf{BTF:    }Set when both the shift register and the data register are empty/full.
    \item \textbf{TxE:    }Set when the data register is empty and ready for new data.
    \item \textbf{RxNE:   }Set when the data register contains received data. 
    \item \textbf{AF:     }This is established when no acknowledgment is received from the slave.
\end{itemize}

\subsubsection{STM32F405 I$^2$C Sequences}


To ensure maximum fidelity between the I$^2$C protocol of the real hardware and 
the I$^2$C of the emulated hardware, not only must the individual logic of each register 
be the same as that of the real hardware, but the behavior of the registers as a whole 
must comply with the expected generation of events of the I$^2$C protocol.

The STM32F405 microcontroller reference manual defines the transmit and 
receive sequences necessary to establish I$^2$C communication. These sequences contain 
the precise operations and register transitions that the firmware must perform at 
each stage of the communication. A model emulating the behavior of this I$^2$C 
interface must be able to recognize these operations and register transitions and 
generate the corresponding I$^2$C events, such as sending the START, STOP, or ACK signals.
These sequence are represented in the Annexe~\ref{ann:i2c_sequences}


%%%%%%%%%%%%%%%%%%%%%%%%%%%%%%%%%%%%%%%%%%%%%%%%%%%%%%%%%%%%%%%%%%%%%%%%%%%%%%%%%%%%%%%%%%%%%%%%%%%%%%%%%%%%%%

%In the Embedded Systems I course of the Master's in Electronic Systems at 
%the UPM, the STM32F4-DISCOVERY development board is used as part of the 
%course materials. This board incorporates an STM32F407VGT6 microcontroller 
%with a 32-bit Arm Cortex-M4 processor. However, since the STM32F407VGT6 
%is not yet directly implemented in the version of QEMU used in SIVALSE, the 
%STM32F405 microcontroller has been selected. It shares many similarities with 
%the STM32F407\cite{STM32F405Datasheet2026}\cite{STM32CortexM42026} and, in addition to the I²C protocol and other peripherals, 
%already has a solid implementation base. Since this thesis focuses on emulating the 
%STM32F405's I²C protocol, this section will list some of its features, taken directly 
%from its datasheet.
%
%The STM32F405RG, part of STMicroelectronics' STM32F4 family, is based on 
%a 32-bit ARM Cortex-M4 core with a Floating Point Unit (FPU), support for DSP 
%instructions, an Adaptive Real-Time Accelerator (ART) for 0-wait execution from 
%Flash up to 168 MHz, a Memory Protection Unit (MPU), and 210 DMIPS of 
%performance. This configuration makes it an ideal microcontroller for embedded 
%systems with real-time and intensive processing requirements, such as motor 
%control, power conversion, and multimedia applications. It incorporates up to 1 MB 
%of Flash memory, up to 192+4 KB of SRAM (with 64 KB of CCM), and essential 
%peripherals: three 12-bit ADCs, two DACs, 17 timers, 16-stream DMA, three I$^2$C ports, 
%multiple USART/SPI/CAN ports, Ethernet MAC, USB OTG HS/FS, and up to 
%140 fast, 5V-tolerant GPIOs. Its memory architecture maps peripherals between 
%0x40000000-0x5FFFFFFFF, as it shows in the table \ref{tab:ARM_memory_map}
%
%\begin{table} [H]
%    \centering
%    \caption
%        {
%            STM32F405 Memory Map
%        }
%    \label{tab:ARM_memory_map}
%\end{table}


%The present thesis is centered in the emulation of the STM32F405's I$^2$C protocol. 
%In this section, will be listed some of the STM32F405 features, extracted directly 
%from its datasheet\cite{STM32F405Datasheet2026}\cite{STM32CortexM42026}. This microcontroller is a high-performance 32-bit 
%microcontroller from STMicroelectronics, built around an Arm Cortex-M4 core with FPU, 
%and DSP extensions, and an adaptive real-time accelerator that enables zero-wait-state 
%code execution from Flash at clock frequencies up to 168 MHz. It delivers up 
%to around 210 DMIPS and includes a memory protection unit, making it suitable for 
%complex real-time embedded applications that require both high performance and 
%robust software isolation.
%
%In terms of memory, the STM32F405 integrates up to 1 MByte of embedded Flash 
%and up to 196 KBytes of SRAM, including a 64 KByte core-coupled memory (CCM) 
%region optimized for fast data access, along with OTP memory and a flexible static 
%memory controller for external SRAM, PSRAM, NOR, NAND, and CompactFlash 
%devices. The device operates from 1.8 V to 3.6 V, offers multiple clock sources 
%(external crystal, internal 16 MHz RC, and 32 kHz oscillators for RTC), and supports 
%several low-power modes (Sleep, Stop, Standby) with a dedicated VBAT domain for 
%RTC and backup registers, enabling energy-efficient designs that retain timekeeping 
%and critical data across power cycles.
%
%The STM32F405 features rich analog capabilities, including three 12-bit ADCs 
%with multi-channel and interleaved modes reaching multi-MSPS aggregate sampling 
%rates, as well as two 12-bit DACs for signal generation. It also provides a 
%16-stream general-purpose DMA controller with FIFO and burst support, a CRC unit, 
%a true random number generator, and a large timer subsystem with up to 17 timers 
%(16- and 32-bit) supporting capture/compare, PWM generation, pulse counting, 
%and quadrature encoder interfaces, making it well suited for motor control, 
%power conversion, and advanced control applications.
%
%On the digital I/O and connectivity side, the STM32F405 exposes up to 140 
%GPIOs (most of them fast and 5 V tolerant) and supports a wide range of communication 
%peripherals: up to three I$^2$C controllers, multiple USART/UART interfaces 
%with support for protocols like LIN, IrDA, and ISO 7816, up to three SPI/I2S blocks, 
%two CAN 2.0B controllers, SDIO, and a parallel camera interface. Advanced connectivity 
%is further enhanced by an integrated Ethernet MAC with IEEE 1588v2 
%hardware timestamping and DMA, as well as full-speed and high-speed USB 2.0 
%OTG controllers with on-chip PHY and dedicated DMA, enabling the microcontroller 
%to serve as a communication hub in networked and multimedia embedded 
%systems.
%
%For development and system integration, the MCU includes a serial-wire/JTAG 
%debug interface, an Embedded Trace Macrocell for instruction tracing, an LCD 
%parallel interface, a real-time clock with calendar and subseconds, and a 96-bit 
%unique device ID that can be used for traceability or security-related features.